\section{Análisis bayesiano}
\label{sec:bayes}

El análisis bayesiano complementa la inferencia frecuentista mediante una
representación probabilística de la incertidumbre sobre los parámetros.
Se desarrollaron dos ejercicios:

\begin{enumerate}
    \item un modelo Beta--Binomial conjugado para la proporción global de
    malos créditos;

    \item una regresión logística bayesiana estimada mediante
    Metropolis--Hastings para estudiar la distribución posterior del
    coeficiente de duración.
\end{enumerate}

El primer ejercicio permite obtener la distribución posterior de forma
analítica y evaluar la sensibilidad de los resultados frente a distintas
distribuciones previas. El segundo requiere simulación MCMC porque la
regresión logística no dispone de una posterior conjugada de forma
cerrada.

\subsection{Modelo Beta--Binomial para la tasa de mal crédito}
\label{subsec:beta-binomial}

En el enfoque frecuentista, la tasa de mal crédito se trató como un
parámetro fijo pero desconocido. A partir de:

\[
x=300,
\qquad
n=1000,
\]

se obtuvo la estimación:

\[
\widehat{p}
=
\frac{x}{n}
=
0.300.
\]

En el enfoque bayesiano, la incertidumbre sobre \(p\) se representa
mediante una distribución de probabilidad. Después de observar los datos,
la distribución previa se actualiza aplicando el teorema de Bayes:

\[
\pi(p\mid x)
\propto
L(x\mid p)\pi(p),
\]

donde:

\begin{itemize}
    \item \(\pi(p)\) es la distribución previa o \textit{prior};

    \item \(L(x\mid p)\) es la verosimilitud aportada por los datos;

    \item \(\pi(p\mid x)\) es la distribución posterior.
\end{itemize}

El resultado no es únicamente una estimación puntual, sino una
distribución completa que describe qué valores de \(p\) son más
compatibles con el prior y la información observada.

\subsubsection{Verosimilitud binomial}

Sea \(X\) el número de malos créditos observados entre \(n\) registros.
El modelo de muestreo es:

\[
X\mid p
\sim
\operatorname{Binomial}(n,p).
\]

La función de probabilidad es:

\[
\Pr(X=x\mid p)
=
\binom{n}{x}
p^x(1-p)^{n-x}.
\]

Para la inferencia sobre \(p\), los términos que no dependen del
parámetro pueden omitirse. La verosimilitud queda expresada como:

\[
L(p\mid x)
\propto
p^x(1-p)^{n-x}.
\]

En el problema analizado:

\[
L(p\mid x)
\propto
p^{300}(1-p)^{700}.
\]

\subsubsection{Distribución previa Beta}

Se asigna a la tasa una distribución previa:

\[
p
\sim
\operatorname{Beta}(\alpha,\beta).
\]

Su densidad es:

\[
\pi(p)
=
\frac{
    1
}{
    B(\alpha,\beta)
}
p^{\alpha-1}
(1-p)^{\beta-1},
\qquad
0<p<1,
\]

donde \(B(\alpha,\beta)\) es la función Beta.

Omitiendo la constante de normalización:

\[
\pi(p)
\propto
p^{\alpha-1}
(1-p)^{\beta-1}.
\]

La media previa es:

\[
\operatorname{E}[p]
=
\frac{\alpha}{\alpha+\beta}.
\]

Los parámetros \(\alpha\) y \(\beta\) controlan tanto la localización como
la concentración de la información previa.

\subsubsection{Posterior conjugada}

Al combinar el prior y la verosimilitud:

\[
\begin{aligned}
\pi(p\mid x)
&\propto
p^x(1-p)^{n-x}
p^{\alpha-1}(1-p)^{\beta-1}\\
&\propto
p^{\alpha+x-1}
(1-p)^{\beta+n-x-1}.
\end{aligned}
\]

Este núcleo corresponde nuevamente a una distribución Beta. Por tanto:

\[
p\mid X=x
\sim
\operatorname{Beta}
\left(
    \alpha+x,\,
    \beta+n-x
\right).
\]

La actualización puede interpretarse como la incorporación de \(x\)
malos créditos al primer parámetro y \(n-x\) buenos créditos al segundo:

\[
\alpha_{\mathrm{post}}
=
\alpha+x,
\]

\[
\beta_{\mathrm{post}}
=
\beta+n-x.
\]

Esta propiedad se denomina conjugación. Permite obtener la posterior
analíticamente, sin recurrir a métodos de simulación.

\subsection{Resúmenes de la distribución posterior}

Para cada prior se calcularon tres resúmenes: la media posterior, la
estimación MAP y el intervalo creíble del \(95\,\%\).

\subsubsection{Media posterior}

La media posterior es:

\[
\operatorname{E}[p\mid x]
=
\frac{
    \alpha+x
}{
    \alpha+\beta+n
}.
\]

Esta cantidad representa el valor esperado de la tasa después de combinar
la información previa con los datos.

\subsubsection{Estimación MAP}

La estimación de máxima densidad posterior es:

\[
\widehat{p}_{\mathrm{MAP}}
=
\frac{
    \alpha+x-1
}{
    \alpha+\beta+n-2
},
\]

siempre que los dos parámetros posteriores sean mayores que uno.

El MAP identifica el valor de \(p\) donde la densidad posterior alcanza
su máximo.

\subsubsection{Intervalo creíble}

El intervalo creíble central del \(95\,\%\) se obtiene con los cuantiles
\(0.025\) y \(0.975\) de la posterior:

\[
\Pr
\left(
    p_{\mathrm{inf}}
    \leq
    p
    \leq
    p_{\mathrm{sup}}
    \mid x
\right)
=
0.95.
\]

A diferencia de un intervalo de confianza frecuentista, este intervalo
permite una interpretación probabilística directa: condicionados al
prior y a los datos observados, la probabilidad posterior de que \(p\)
se encuentre entre ambos límites es del \(95\,\%\).

\subsection{Distribuciones previas evaluadas}

Se compararon tres priors con distintos niveles de información.

\subsubsection{Prior uniforme}

El primer modelo utiliza:

\[
p
\sim
\operatorname{Beta}(1,1).
\]

Esta distribución es uniforme en el intervalo \([0,1]\):

\[
\pi(p)=1.
\]

Antes de observar los datos, asigna la misma densidad a cualquier valor
posible de la tasa. Se utiliza como una especificación inicial poco
informativa.

La posterior resultante es:

\[
p\mid x
\sim
\operatorname{Beta}(301,701).
\]

\subsubsection{Prior de Jeffreys}

El segundo modelo utiliza:

\[
p
\sim
\operatorname{Beta}(0.5,0.5).
\]

El prior de Jeffreys es una alternativa no informativa construida a
partir de la información de Fisher del modelo binomial. Su densidad
asigna mayor peso a los extremos del intervalo que el prior uniforme.

La posterior es:

\[
p\mid x
\sim
\operatorname{Beta}(300.5,700.5).
\]

Su inclusión permite comprobar si una definición no informativa diferente
produce cambios apreciables en la inferencia.

\subsubsection{Prior informativo débil}

El tercer modelo utiliza:

\[
p
\sim
\operatorname{Beta}(6,18).
\]

Su media previa es:

\[
\operatorname{E}[p]
=
\frac{6}{6+18}
=
0.25.
\]

La concentración total es:

\[
\alpha+\beta
=
24.
\]

Este prior representa una expectativa inicial de que la tasa se encuentre
alrededor del \(25\,\%\), pero con una concentración pequeña frente a las
\(1000\) observaciones disponibles. Por ello, se considera informativo
débil.

La posterior resultante es:

\[
p\mid x
\sim
\operatorname{Beta}(306,718).
\]

El propósito de este prior no es imponer una tasa de \(25\,\%\), sino
evaluar cuánto se desplaza la inferencia cuando se incorpora una
expectativa previa moderadamente distinta de la tasa observada.

\subsection{Resultados posteriores}

La Tabla~\ref{tab:resultados-beta-binomial} presenta los resultados para
los tres priors.

\begin{table}[H]
    \centering
    \caption{Resultados del modelo Beta--Binomial bajo diferentes
    distribuciones previas.}
    \label{tab:resultados-beta-binomial}

    \small
    \begin{tabularx}{0.98\textwidth}{
        @{}
        >{\raggedright\arraybackslash}p{3.4cm}
        >{\centering\arraybackslash}p{2.5cm}
        >{\centering\arraybackslash}p{1.6cm}
        >{\centering\arraybackslash}p{1.8cm}
        X
        @{}
    }
        \toprule
        Prior
        & Posterior
        & MAP
        & Media
        & Intervalo creíble del \(95\,\%\) \\
        \midrule

        \(\operatorname{Beta}(1,1)\), uniforme
        & \(\operatorname{Beta}(301,701)\)
        & \num{0.3000}
        & \num{0.3004}
        & \([0.2724,\;0.3291]\) \\

        \(\operatorname{Beta}(0.5,0.5)\), Jeffreys
        & \(\operatorname{Beta}(300.5,700.5)\)
        & \num{0.2998}
        & \num{0.3002}
        & \([0.2722,\;0.3289]\) \\

        \(\operatorname{Beta}(6,18)\), informativo débil
        & \(\operatorname{Beta}(306,718)\)
        & \num{0.2984}
        & \num{0.2988}
        & \([0.2712,\;0.3272]\) \\

        \bottomrule
    \end{tabularx}
\end{table}

Los tres modelos producen resultados muy próximos. Las estimaciones MAP
se encuentran entre:

\[
0.2984
\quad\text{y}\quad
0.3000,
\]

mientras que las medias posteriores se sitúan entre:

\[
0.2988
\quad\text{y}\quad
0.3004.
\]

El prior uniforme produce un MAP de \(0.3000\), coincidente con la
estimación frecuentista. El prior de Jeffreys genera un resultado casi
idéntico.

El prior \(\operatorname{Beta}(6,18)\) desplaza ligeramente la posterior
hacia su media previa de \(0.25\). Sin embargo, la media posterior solo
disminuye hasta \(0.2988\), lo que muestra que la influencia del prior es
reducida frente al tamaño de la muestra.

\subsection{Sensibilidad al prior}

La Figura~\ref{fig:beta-binomial} compara las distribuciones previas y
posteriores.

\begin{figure}[H]
    \centering

    \includegraphics[
        width=\textwidth
    ]{
        figura_10_modelo_beta_binomial.png
    }

    \caption{
        Sensibilidad del modelo Beta--Binomial a la elección del prior.
        El panel izquierdo muestra las tres distribuciones previas. El
        panel derecho presenta las posteriores, la estimación
        frecuentista y el intervalo de Wilson del \(95\,\%\).
    }

    \label{fig:beta-binomial}
\end{figure}

En el panel izquierdo, los priors presentan formas claramente diferentes.
El uniforme asigna densidad constante, Jeffreys concentra más densidad
cerca de cero y uno, y el prior informativo concentra su masa alrededor
de valores próximos a \(0.25\).

Después de observar los \(1000\) registros, las tres posteriores se
superponen casi completamente. La evidencia contenida en los datos domina
las diferencias entre las distribuciones previas.

La amplitud de los intervalos creíbles también es muy similar. Entre el
límite inferior más bajo y el más alto existe una diferencia de apenas:

\[
0.2724-0.2712
=
0.0012,
\]

equivalente a \(0.12\) puntos porcentuales. Entre los límites superiores,
la diferencia máxima es:

\[
0.3291-0.3272
=
0.0019,
\]

equivalente a \(0.19\) puntos porcentuales.

Estas variaciones no modifican la conclusión sustantiva: la tasa de mal
crédito se concentra alrededor del \(30\,\%\).

\begin{hallazgo}
La inferencia sobre la tasa de mal crédito es robusta frente a los tres
priors considerados. Con \(n=1000\), los datos dominan la actualización
y mantienen la media posterior próxima a \(0.30\).
\end{hallazgo}

\subsection{Comparación con la inferencia frecuentista}

En la Sección~\ref{sec:incertidumbre}, el intervalo de Wilson fue:

\[
\operatorname{IC}_{95\,\%,\,\mathrm{Wilson}}
=
[0.2724,\;0.3291].
\]

Con el prior uniforme, el intervalo creíble fue:

\[
\operatorname{ICr}_{95\,\%,\,\mathrm{Beta}(1,1)}
=
[0.2724,\;0.3291].
\]

Los límites coinciden numéricamente al nivel de redondeo utilizado. Esta
coincidencia no significa que ambos intervalos sean conceptualmente
iguales.

El intervalo de Wilson se interpreta como una propiedad de un
procedimiento repetido: en un conjunto hipotético de muestras comparables,
aproximadamente el \(95\,\%\) de los intervalos construidos mediante
Wilson contendría el verdadero valor de \(p\).

El intervalo creíble responde a una afirmación distinta:

\[
\Pr
\left(
    0.2724
    \leq
    p
    \leq
    0.3291
    \mid x
\right)
=
0.95,
\]

condicionada al modelo binomial, al prior uniforme y a los datos
observados.

La similitud numérica se explica por el tamaño de la muestra y por el uso
de un prior débil. En estas condiciones, tanto la aproximación
frecuentista como la posterior se concentran cerca de la proporción
muestral.

\begin{table}[H]
    \centering
    \caption{Comparación entre la inferencia frecuentista y bayesiana
    para la tasa de mal crédito.}
    \label{tab:comparacion-frecuentista-bayes}

    \begin{tabularx}{0.94\textwidth}{
        @{}
        >{\raggedright\arraybackslash}p{3.5cm}
        >{\centering\arraybackslash}p{2.2cm}
        >{\centering\arraybackslash}p{3.0cm}
        X
        @{}
    }
        \toprule
        Enfoque
        & Estimación
        & Intervalo del \(95\,\%\)
        & Interpretación \\
        \midrule

        Frecuentista: Wilson
        & \(\widehat{p}=0.3000\)
        & \([0.2724,\;0.3291]\)
        & Cobertura del procedimiento bajo muestreo repetido. \\

        Bayesiano:
        \(\operatorname{Beta}(1,1)\)
        & MAP \(=0.3000\)
        & \([0.2724,\;0.3291]\)
        & Probabilidad posterior condicionada al prior y los datos. \\

        \bottomrule
    \end{tabularx}
\end{table}

Ambos enfoques conducen a la misma conclusión aplicada: la tasa de mal
crédito del conjunto analizado se estima cerca del \(30\,\%\), con una
incertidumbre aproximada de tres puntos porcentuales a cada lado.

\subsection{Supuestos y alcance del modelo}

El modelo Beta--Binomial supone que, condicionado a \(p\), los registros
pueden representarse como ensayos Bernoulli comparables e independientes.
También supone una única tasa común para toda la muestra.

Esta simplificación ignora posibles diferencias entre propósitos,
duraciones, perfiles financieros o periodos de originación. El modelo
describe la tasa agregada de la cartera, pero no sustituye la regresión
logística utilizada para estimar probabilidades condicionales.

La posterior tampoco corrige problemas de representatividad de la
muestra. Una distribución posterior estrecha puede reflejar alta precisión
condicionada al dataset y, al mismo tiempo, coexistir con una validez
externa limitada.

Por tanto, el intervalo creíble cuantifica la incertidumbre estadística
sobre la tasa dentro del modelo asumido. No garantiza que la misma tasa
sea aplicable a otra institución, país o periodo.

% ============================================================
% En el siguiente bloque se incorporarán:
% - Regresión logística bayesiana
% - Prior normal débilmente informativo
% - Metropolis--Hastings multivariante
% - Diagnósticos de la cadena
% - Posterior del coeficiente de duración
% ============================================================

% Contenido previsto:
% - Modelo Beta--Binomial.
% - Prior, verosimilitud y posterior conjugada.
% - Media posterior, MAP e intervalo creíble al 95 %.
% - Sensibilidad a por lo menos dos priors.
% - Comparación con el IC de Wilson.
% - Metropolis--Hastings para un coeficiente logístico.
% - Diagnóstico de cadena e interpretación posterior.
